\section{The Foundation of Rome}

The racial origins of Rome are surrounded by the mystery that accompanies the birth of every great civilization: fable and history intermix with each other. Beyond the veil of myth, the foundation of the City of Rome presents in its drama the signs of a new history of the West, through a chain of events that are to be considered symbols of a new way of seeing and organizing life and that therefore, in the old and decadent Mediterranean world, herald the Roman conquest of the West. It should be said that the rite of the foundation is the prelude, in its naked drama, to the eternal virtues of the Empire: the founder is in fact the \emph{Lar} par excellence, the progenitor of a race of divine origin, and first and foremost achieves a spiritual pact with the gods and with men; then he will proceed to the rite of the foundation.

In order to understand the meaning of a power that initiates the formation of the race, it is necessary to keep in mind that everything in Rome comes from a sacred, supernal sense of life: the foundation of the city is a constructive act that moves from an order of purely spiritual necessity: it is the result of a religious ideal harmony constituted among a group of men, warriors and priests, who seek to establish a new unity, a creative force, and intend primarily to locate the sanctuary of common worship in the city. The foundation, then, is not motivated by economic, topographic, or, in any case, contingent reasons, but by reasons of common aspiration to a higher ideal of a supermaterial, ``divine'' level. The rite then discloses a universalistic value assigning celebrations of nature, mystical and esoteric meanings, queries to the celestial forces for the choice of location and evocations to demonic powers, the genius of the people, to the chthonic and uranic gods, for the fortune of the city that is going to arise.

While wanting to mention this \emph{mystery of the founding of Rome} in this paper, it is not possible to linger at this point to give even a brief idea of the essential value attached by the ancients, and especially by the Romans, to the action of the rite. Suffice it to say that, in the same way that a modern with mechanical operations and procedures becomes master of distance, it gives form to matter and organizes its own outer life, so the ancient Roman Empire, through the ritual technique made perfect thanks to the regal-priesthood union that entailed the action of a self-conscious, ``solar'' will, and the body of a dynamic, mediating, ``lunar'' force, established an ascending contact with the magnetic forces of the universe and acted psychically through them. There existed a science of such action: unlike the mechanical sciences placing all men on the same levels (since mechanical means can be operated both by the learned and the ignorant), it required a spiritual dignity that was not for everyone; it demanded the presence of psychical qualities, in a dynamic and exceptional sense, and hence connected to a higher morality that was not at all different from that of the mystic, the priest. However, this does not mean for the Roman that the event must be limited to the contemplative and mystical world (as in the ancient rituals of the Mediterranean peoples, Orphism and Pythagorism) but from a spiritually ``superhuman'' plane, it was necessary to speak in order to give meaning to the ``human'', to the real, to everyday life, to organic political unity. It was therefore a highly religious sense of that existence which the rite conformed to: the race was continuously moulded in it, made strong in every aspect: and they were privileged beings, that is, more inwardly complex, far-sighted, ``Initiated'', and were entrusted with the task of giving shape and direction to events, through the meticulous technique of the rite. They were priests, flamens, pontiffs. Their wisdom was the Tradition, the very remote legacy of the ``solar'' race, heroic and spiritual at the same time.

With the founding of Rome, the priesthood resumed to be part of an eternal harmony of the world, becoming a complement of warrior action, that is, no longer limited to the aestheticism of the ancient Western mysteriological communities (Zagreus, Dionysus, Orpheus), nor serving as an instrument for abusive orgies and to matriarchal ``aphroditism'' of previous Mediterranean regimes, but compelling the Spirit in a new discipline, that is, to result in precision of acts, assertions, conquests, constraining the Spirit in a new discipline, that is, to result in a precision of acts, affirmations, conquests, obliging the \emph{apeiron}, the indefinite, to assume form, to be actualized in \emph{pèras}, in well-defined reality, the archetype of the classical creation and the original inspiration of the \emph{imperium}, of the new architectonic soul.

We repeat, this is not poetry, nor even myth. Despite the magical resonance of myth in the story of Aeneas, as Virgil relates it, and in that of Romulus, son of a Vestal Virgin and the god Mars, although later than the events of the origins, it represent only a human veil of deep mystery of the foundation. Someone told us that with this legend we tried to make Romulus an emanation of the ``feminine'' principle enlivened by the Numidian force personified in the masculine God. Those, who believe that it's worth the trouble to attribute value to the analogies, keep in mind that she is to be considered as the terrestrial symbol of virginity despite the union of Mars with the Vestal Virgin, that is, of the raw power not actualized in motherhood.

The essence of the secret can be glimpsed only if, taking advantage of a non-scholastic and non-rationalistic vision of history, we consider that Romulus grafted onto it acts which have new meanings, while adopting the archaic rite of the foundation. It is not enough to recognize that this rite, although of Etruscan origin, was also common to Lazio and Sabina. In the rite of the \emph{mundus} the principle of the \emph{eternity} of the Urbe is realized, since the Spirit again results in action, in hierarchical reality. For those interested in it, we refer to symbols, like the prophetic visions of Romulus on the Palatine, and emphasize that Remus, who is to symbolize the anti-hierarchical element of the decadent period, breached the intangibility of the trench and Romulus punished him.

According to the archaic Etruscan rite, the auguries had to arise after midnight, in silence, and wait for the dawn. Romulus and Remus also rise \emph{post mediam noctem}: they go up onto two hills (\emph{tabernacula capiunt}, \emph{templa capiunt}): from this moment the destiny of Rome and of its race is going to be marked. Historians and poets almost totally agree in telling us that Romulus went up on the Palatine and Remus on the Aventine: two different places, two opposed symbols, two traditions that clash, and hence still two races.

It is necessary to decide on the name of the new city: will it be called Roma or Remora? Will the king be Romulus or Remus? All are intent, in anticipation of the verdict of which must come from the same force of destiny. The white disk of the moon sets: it diffuses the light of dawn and here is the most perfect of the auguries: the eagle of Jupiter shows up on the left --- it is the symbol of ``Olympian'' royalty characteristic of the ``Solar Race'', which is revealed to the fathers of the future rulers of the West --- and while the disc of the sun appears, there is a flock of black birds flying rapidly. The one who had seen the first twelve vultures would reign. First is Romulus, at the dawn of the day; the people rejoice: Romulus is consecrated king and priest, the first \emph{Lar}, the father of the new race.

And so it is that an authentic progenitor demonstrates the priestly technique of the foundation. He, informed of the ancient Etruscan ritual learned through the secret liturgical books --- as we read in Cato, Servius, Festus and Gellius --- initiated a very spiritual sacred science which integrated in him the warrior and the founder of civilization, he drew the omens, offered the sacrifice, lighted the ritual fire, dug the circular pit, the \emph{mundus}, and threw on it the handful of earth to which the soul of ancestors was symbolically and actually tied, he began the powerful and mysterious life of the \emph{terra patrum}, the land of the fathers, the fatherland, that is, of the earth to which the destiny of the race will be linked.

To seal the bond of the needed deity with the spatial center of the new city, that is, in order \emph{to link the power of the Spirit to the place}, so the place contains \emph{his} ``demonic'' forces of the fatherland, the sacred place, the actual eternity, \emph{a large stone}, the \emph{lapis Manalis}, \emph{closes the mouth of the pit}.

The ``underworld'' is so constituted, that it must accept the souls, not the bodies, of the dead, and hence three times a year they will issue in the world of life. Delivered to the underworld, there are erected a conically shaped column and a pyramid: both are sacred to the \emph{manes} of the progenitors and are consecrated to his heroism. It is therefore an immortal strength that goes to the ground, which will therefore also be immortal. After taking in the divine cycle, the founder, living spiritually in the underworld, will be revered by the city as a son of the gods, a god among gods, \emph{auctor}, hero and kin of the new people.

With the underworld and the supernal word consecrated, we proceed to the ritual constitution of the topography of the city, always with regard to an ancient ceremonial secret that Romulus knows well. We know only the exterior modality of the ceremonial, but it also has a language for those who can understand it. In a white cloak and his head covered, according to the priestly custom, he yoked an ox and a white robust cow to the plow, he descended down the hill, followed by silent companions, and invoking with mysterious formulas of propitiation the favor of divine forces, began to trace the ritual trench, making sure that inside, on the side of the city, is the cow, the image of fertility, and outside, on the countryside, is the ox, a symbol of strength.

In driving the trench, he lifts the plow where he wants the gates, so that it does not touch the ground. Then he raises the walls of the city limits, following the line of the trench, and outside, hugging the walls, dig the trench of the circumvallation: here and there the two \emph{pomeri}: one inside and the other outside: two land areas that they can neither plow nor inhabit, purposely vacated and free, for the purpose of lookout and defense. The sacred walls will be raised here and no one will be able to then modify the size and restore them without the permission of the Pontiffs. At the borders the names dedicated to the god Terminus will be placed.

Having traced the borders of the city, he gave houses to the fathers subject to chance, declared the laws, and, followed by all the comrades, climbed back up the summit. Then, he shouted out the divine name of the city, which is repeated three times in a loud voice by the fathers, immolated the white bull with the cow on the altar of the Supreme Jupiter. They then prepared meals and celebrations lasting for nine days. The objects used in the rite of the foundation of the Urbe were put back into the \emph{mundus} as sacred.

This complex ritual so that Rome, according to Ennius, was founded with an ``august omen,'' contains the basic reasons that will give the sense of eternity to the race, the city and its empire: it is the ceremonial aspect of a secret technique aspiring to subjugate the events to a single direction, that of the nascent Urbe . It is the initial victory of the race of Rome over fate, for a new cycle of the West. This will be then the meaning of the \emph{Dies Natalis Urbis Romae}. The founding of Rome is thus a constructive act that moves from an order of inner necessity: while it is the consequence of a religious agreement among those who will live there, since it will represent the sanctuary of common worship, it must be ritually established as \emph{cause of causes} , as a starting point, as a radical design of a future organism. It is a seed in the womb of the earth and, as a seed, it must contain the force of procreation.

It is essentially an ``initiatic'' art that intervenes to give direction to \emph{destiny}, with the rite of the \emph{mundus}: this small circular pit, dug by Romulus, accepts a handful of earth that he has brought with him from Alba (Plutarch, Dione, Cassius, Ovid, Festus) and welcomes the sod that each of his companions has carried from the land in which it was burning before the sacred fire, which was linked to the souls of their \emph{manes}. It is therefore land steeped in forces, the soul of races, the land where the dynamics of the \emph{genius loci}, the spirit of the race is attached. It is not poetry. It is the creation of a powerful condenser of forces gathered according to a process whose modality is not known by many, and indeed must be unknown, --- always in compliance with an esoteric technique --- just as the secret name of the city, the \emph{nomen sacrum}, the seminal word, \emph{spermaticos logos}, the secret word that corresponds to the virtues of the deity of the city. The same force of the rite ensures that the founders are liberated from the ``impiety'' of leaving the land of their forefathers and promotes a new link between soul and matter, between the spirit of the race and the land chosen for the foundation.

The \emph{mundus} is, therefore, a sacred place, the central point in which destiny is tied to the earth through ritual strength: the spatial aspect of becoming is thus dominated and bound by virtue of an incident that contains in itself the surpassing of time: if the land is linked to the power of the Spirit, if the Spirit of the deified ancestor is immortal, the land is saturated in a metaphysical virtue, it becomes the mystical center of eternity. \emph{Mundus} means, in ancient esoteric language, the region of the Manes, \emph{moundos} (Plutarch, Festus, Servius); and because the worship of the Manes is uninterrupted due to burning the sacred flame in front of the domestic \emph{lararium}, it also clarifies the meaning of the \emph{fire}. The sacred fire of the city will be lighted on the prophetic pit, which will be the \emph{fire} of Vesta, always lit in the temple: it will not be an element of deified nature, as historical criticism has always believed, but will represent the earthly symbol of a divine force that, on the celestial plane, will always symbolically correspond to the sun and in the sense of human physiology, the heart, the seat of superhuman intelligence according to ancient spirituality (Cicero, Plotinus, Iamblichus, Emperor Julian). Thus as an eternal flame of divinity burns in the heart of the hero and the ascetic, the fire of Vesta will burn inside of the temple.

But who in Rome first lit this fire? Romulus. He then is the founder, but also the primigenial \emph{Lar} of the city, the spiritual father of the Roman race: and being divine in human life, his death will only be a total reuniting with the divine plane. The myth dramatizes this story.

What is necessary to emphasize is that the divine element constitutes the essential part of the birth of Rome, it is only the religious aspect of the mastery of events, of inevitability, obtained through the possession of transcendent energy that was familiar to the ancient Initiate, as the control and mastery of physical energies is to the contemporary engineer and mechanic.

\begin{quotex}
From \textit{La Razza di Roma}\footnote{\url{https://www.gornahoor.net/library/LaRazzaDiRoma.pdf}}, by \textbf{Massimo Scaligero} (1939)
\end{quotex}

\flrightit{Posted on 2021-10-05 by Aeneas}

